% ==========================================
% BÖLÜM 2: KULLANICI İÇİN TAVSİYELER
% ==========================================
\chapter{Kullanıcı İçin Tavsiyeler}

Sistemin kararlı, hızlı ve veritabanı bütünlüğünü bozmadan çalışabilmesi için son kullanıcıların dikkat etmesi gereken temel operasyonel hususlar aşağıda sunulmuştur:

\vspace{0.3cm}

\begin{tcolorbox}[colback=lightgray,colframe=bordergray,leftrule=3.5pt,boxrule=0.4pt,top=6pt,bottom=6pt]
    \textbf{\textcolor{primary}{\faMousePointer\ Sayfa Hücrelerini Manuel Değiştirmeyin:}}\\
    \small Uygulama, arka planda dinamik koordinatlama ve kesişim tespiti yapmaktadır. \texttt{Yenile} veya veri sayfalarındaki satır/sütun yapılarını elle silmek veya yerlerini değiştirmek arka plan servislerinin (\texttt{mod\_Services}) yanlış hücrelere yazmasına yol açabilir.
\end{tcolorbox}

\vspace{0.2cm}

\begin{tcolorbox}[colback=lightgray,colframe=bordergray,leftrule=3.5pt,boxrule=0.4pt,top=6pt,bottom=6pt]
    \textbf{\textcolor{primary}{\faSave\ Veri Kaydetme Adımlarını Sırayla Takip Edin:}}\\
    \small Paneller üzerindeki \textbf{"Güncelle"} butonlarına basılmadan \textbf{"İleri"} adımlarına geçilmemelidir. Butonlar yeşil renge dönüştüğünde verinin ilgili veritabanı hücresine güvenle işlendiği teyit edilmiş olur.
\end{tcolorbox}

\vspace{0.2cm}

\begin{tcolorbox}[colback=lightgray,colframe=bordergray,leftrule=3.5pt,boxrule=0.4pt,top=6pt,bottom=6pt]
    \textbf{\textcolor{primary}{\faCalendarDay\ Geçmiş Ay Taramalarında Format Uyumluluğu:}}\\
    \small Geçmiş ay verilerini sorgularken açılır menüden doğru yıl ve ay seçiminin yapıldığından emin olunuz. Sayfa isimlerinin ve tarih formatlarının sistem standartlarına (\texttt{dd.mm.yyyy}) uygun olması sorgu hızını doğrudan artıracaktır.
\end{tcolorbox}

\vspace{0.2cm}

\begin{tcolorbox}[colback=yellow!15,colframe=orange!80!black,leftrule=3.5pt,boxrule=0.4pt,top=6pt,bottom=6pt]
    \textbf{\textcolor{orange!80!black}{\faLightbulb\ Altın Kural - Düzenli Yedekleme:}}\\
    \small Toplu güncelleme veya veri aktarım operasyonları öncesinde çalışma kitabının bir yedeğinin alınması, olası kullanıcı hatalarında veri kaybını önlemek adına tavsiye edilir.
\end{tcolorbox}